
\subsection{User Characteristic}

There are three types of users: Students, Company Members, and University Members. Each user is identified by an email and a password, and once registered, 
they can log in using only these two credentials.

\subsubsection{Student}

The student registers by providing their name, surname, an email associated with the university they belong to (registered on the platform), and a password. The student joins the Students\&Companies platform to find an internship that fulfills their needs. 
They can apply to multiple internship proposals, and if accepted, they can track their experiences through their account on the platform. Each student account is linked to 
the same university associated with their email and is responsible for handling their complaints.

\subsubsection{Company Member}

The company member is a user who represents an employee of a company registered on the Students\&Companies platform. The set of all company members constitutes the company on the system. 
They register by providing their name, surname, an email from their company's domain, and a password. A company member can publish internships offered by their company, 
review student applications, communicate with applicants to arrange interviews. After the interviews are concluded, company members will select the candidates who will be 
hired for the internship.

\subsubsection{University Member}

The university member is a user who represents a worker at a university registered on the Students\&Companies platform. A university member handles all student 
complaints and can interrupt an internship when necessary. The university is represented by the set of all its members.


%todo: specify somewhere that if the university it's not registered the student can't register to the platform (same for university members and company member but with company)
%todo: if student change university he creates a new account
%todo: when a university email expires the student account is deleted by the platform (when periodically or when he tries to log in)

\subsection{Product perspective}
\subsubsection{Scenarios}
\begin{itemize}
      \item Company announces a new internship
            Code S.P.A announces a new software engineering internship in Italy.
            The company's secretary logs into the Student \& Company platform and provides a brief description of the intern's tasks and requirements.
            The system detects that the language requirement is missing and prompts the secretary with the message: "No language requirement?".
            Acting on this suggestion, the secretary adds proficiency in English as a standard requirement.
            Once finalized, the secretary submits the description to the platform.

      \item Marco subscribes to the platform
            Marco creates an account on the platform and prepares his CV using the guidelines displayed on-screen.
            As a master's student in computer science, Marco marks his profile as open to software engineering internship opportunities.

      \item Match accepted
            Marco logs into the platform, eager to secure an internship.
            On the homepage, he is notified of an available software engineering position.
            Intrigued, Marco accepts the match.
            The system sends Marco's application to Code S.P.A, where the recruiter, Anna, reviews it.

      \item Start of the selection process
            The next day, Anna logs into the platform and reviews Marco's profile.
            Impressed by his qualifications, she adds Marco to the list of shortlisted candidates.
            Marco is notified of his selection through a private message from Anna.
            They use the platform's messaging feature to agree on the interview's date and format, ultimately deciding on an online interview.
            Anna schedules the meeting and generates a link to an external videoconference using the platform's built-in tools.

      \item New interns are notified
            Once all interviews are completed, the recruiter updates the system with the final decisions.
            The platform notifies Marco that he has been accepted as an intern.
            Marco logs in to confirm the good news.
      %TODO: add scenario related to the diary
      \item Student writes a complaint
            After starting his internship, Marco finds that the tasks assigned to him do not align with the job description and are expected to be completed within an unreasonably short timeframe.
            Frustrated, Marco uses the platform to write a formal complaint intended for his university.

      \item Company writes a complaint
            
            A student named Pino fails to respond to emails from colleagues and neglects his assigned tasks.
            As a result, the company files a complaint against him through the platform.

      \item University monitors ongoing internships. 
            Zeno, the university's internship coordinator, reviews Marco's complaint and investigates the issue.
            He notices that similar complaints have been made by other students in the past.
            Zeno contacts Daniele, the company's internship manager, to address the issue.
            Daniele assures Zeno that the company will strive to align tasks more closely with the descriptions provided.

      \item Internship is interrupted
            Zeno also reviews the complaints filed against Pino.
            Noting that Pino has repeatedly violated the legal agreement governing his internship, Zeno decides to terminate the internship after the third complaint.
            Through the platform, Zeno informs Code S.P.A of his decision and apologizes for the student's behavior.
            The company reopens the internship position in response.

      \item End of the internship
            At the conclusion of the internship, both Marco and his supervisor provide feedback about their experience.
            This feedback is submitted through the platform and used to enhance the platform's matchmaking algorithms for future internships.

\end{itemize}

\subsubsection{Further details on the shared phenomena}
\subsubsection{Class diagrams}

%TODO: put diagram

The previous diagram aims to model all relevant aspects of the system mentioned so far (e.g., in the phenomena).
Cardinalities of type "1" have been omitted to improve the readability of the diagram.

The first key concept is the user.
\textbf{User} class represents a generic and abstract physical actor interacting with the system.
A user account (email and password) is associated with it, enabling authentication in the system.
There are three types of users: \textbf{Student}, \textbf{CompanyMember}, and \textbf{UniversityMember}. Each "member" is then associated with a
specific institution (\textbf{University} or \textbf{Company}, depending on the user type).

The \textbf{Internship} class represents an internship proposal created by a company.
A few clarifications about its attributes: "maxStudentsAllowed" is the number of students the company aims to recruit
and "openForApplication" specifies whether applications are currently allowed. Each internship proposal can involve multiple students.
For this reason, the \textbf{InternshipStatus} class is introduced. It represents an instance of an internship for a specific student.
This class contains all the information regarding the internship's progress, ensuring that each participant (student, company, and university) stays informed about its status.

There are three classes associated to the InternshipStatus class: InternshipComplaint, Feedback, and InternshipReport.

\textbf{InternshipComplaint} models a complaint created by a student and directed to the company.
\textbf{Feedback} represents the final evaluation written by the student and/or the company about the other party at the end of the internship.
The "globalScore" field allows for an overall score, while the \textbf{FieldEvaluation} class enables detailed feedback on specific aspects. 
These evaluations are used by the system to generate recommendations.
\textbf{InternshipReport} represents a report written by a student or the company. A report is a brief description of what occurred during the internship on a specific day.

The \textbf{Chat} class describes communication between the student and the company member. It is instantiated only after the match is confirmed by both the two parties.

The \textbf{StudentSelectionProcess} class handles all stages of selecting an \textbf{individual} student, from matching to communicating the results (including interviews). The outcome referred
to a specific student is represented by the attribute "selected".

Additionally, the \textbf{Request} class models a "match request" made by either the student or the company and directed to the other party.
If both parties accept, the attribute "approved" will be set to true and then the selection process begins.

Finally, the \textbf{Suggestion} class represents tips that the system provides, either to students (e.g., for their CV) or to companies (e.g., for internship descriptions).

\subsubsection{State diagrams}

\subsection{User characteristics}
\subsection{Assumptions, dependencies and constraints}
\begin{itemize}
    \item \text{[DA1]} The statistical analysis tool produces reliable analysis.
    \item \text{[DA2]} The student writes only correct information inside the CV.
    \item \text{[DA3]} The university correctly interrupts internships, only if it is strictly necessary.
    \item \text{[DA4]} The video conference tool correctly creates the room for interviewer and student.
    \item \text{[DA5]} The company properly writes the description (including terms and requirements) of the internship.
    \item \text{[DA6]} The interviewer correctly conducts the interview.
    \item \text{[DA7]} The company selects students following proper criteria: company doesn't discriminate candidates based on their age, gender, ethnicity and selects the most deserving candidates.
    \item \text{[DA8]} The company schedules the interview based on the time and date agreed with the student.
    \item \text{[DA9]} An interviewer from the company attend to the scheduled interview.
    \item \text{[DA10]} The university consistently provides a precise and comprehensive list of all email addresses assigned to its enrolled students. %specify in separated section
    \item \text{[DA11]} During the registration process, the student must enter his correct university email contact.
\end{itemize}
